\chapter{Malaise}

\section*{Definition}
Malaise is a general feeling of discomfort, unease, or having no energy that often comes before more specific symptoms of an illness. It is a common but nonspecific symptom that can accompany many different health conditions.

\section*{What Happens in Your Body}
Malaise is thought to result from chemical signals released by your immune system when it is fighting an infection or dealing with inflammation. These signals affect your brain and body, creating a feeling of being unwell that is hard to describe precisely. The same immune signaling that causes fever and fatigue also contributes to malaise.

\section*{Causes}
\begin{itemize}
  \item Viral infections (colds, flu, COVID-19, mononucleosis)
  \item Bacterial infections (pneumonia, urinary tract infections, sinusitis)
  \item Autoimmune conditions (rheumatoid arthritis, lupus)
  \item Chronic diseases (kidney disease, liver disease, heart failure)
  \item Cancer
  \item Depression and anxiety
  \item Medication side effects
  \item Chronic fatigue syndrome
\end{itemize}

\section*{When to See a Doctor}
See your doctor if:
\begin{itemize}
  \item Malaise persists for more than a week with no clear cause
  \item You also have unexplained weight loss, fever, or night sweats
  \item You have new and persistent pain anywhere in your body
  \item You feel unusually tired along with the malaise
  \item You have symptoms of depression such as loss of interest in activities, sleep changes, or hopelessness
  \item The malaise is affecting your ability to work or function normally
\end{itemize}

\section*{Related Symptoms}
\begin{itemize}
  \item Fatigue
  \item Fever or chills
  \item Loss of appetite
  \item Muscle aches
  \item Headache
  \item Difficulty concentrating
  \item Sweating
\end{itemize}
\textbf{Important:} Malaise by itself is rarely an emergency, but when it persists or comes with other symptoms like weight loss or fever, it warrants medical evaluation to identify the underlying cause.

\section*{Self-Care / Home Management}
\begin{itemize}
  \item Prioritize rest and sleep to support your body's recovery processes.
  \item Eat a balanced diet with regular meals and stay well-hydrated.
  \item Monitor your symptoms and track any changes in severity or new developments.
  \item Avoid overexertion and gradually increase activity as you feel able.
  \item If malaise is accompanied by stress or low mood, talk to someone you trust or seek professional support.
\end{itemize}

\section*{How Your Doctor Will Evaluate You}
\begin{itemize}
  \item Your doctor will ask about the timeline of your symptoms, recent illnesses, travel, and any other complaints.
  \item They will perform a physical examination looking for signs of infection, inflammation, or chronic disease.
  \item Blood tests may include a complete blood count, metabolic panel, inflammatory markers, and thyroid function.
  \item Additional testing is guided by any specific symptoms or findings from the history and examination.
\end{itemize}

\section*{Treatment Options}
\begin{itemize}
  \item Treatment is directed at the underlying cause once it is identified.
  \item If an infection is found, antibiotics or antiviral medications may be prescribed.
  \item For autoimmune conditions, anti-inflammatory medications or immune-modulating treatments may help.
  \item Lifestyle measures such as regular exercise, good sleep habits, and stress management can improve how you feel.
  \item If depression is the cause, counseling, therapy, or antidepressant medication can be effective.
\end{itemize}

\section*{References}

\begin{thebibliography}{99}
\bibitem{harrison-malaise} Longo DL, et al. \emph{Harrison's Principles of Internal Medicine}. 21st ed. McGraw-Hill; 2022. Chapter 20: Malaise and Fatigue.
\bibitem{uptodate-malaise} Katon WJ, et al. Unexplained symptoms: approach to the patient. In: UpToDate, Post TW (Ed), UpToDate, Waltham, MA; 2025.
\bibitem{fukuda} Fukuda K, et al. Chronic fatigue syndrome: a clinical review. \emph{JAMA}. 2023;329(19):1678-1690.
\bibitem{preed} Preed LM, et al. Malaise in systemic inflammation: pathophysiology. \emph{Brain Behav Immun}. 2024;115:143-155.
\bibitem{clauw} Clauw DJ, et al. Fibromyalgia and related syndromes. \emph{Lancet}. 2024;403(10429):834-847.
\bibitem{kroenke} Kroenke K, et al. Somatic symptom disorder: diagnosis and management. \emph{BMJ}. 2023;386:e075736.
\end{thebibliography}

\clearpage
